\documentclass[10pt,a4paper]{amsart}
\usepackage[margin=19mm]{geometry}
\usepackage{amsmath,amssymb,mathtools}
\usepackage[scr=rsfs,cal=euler]{mathalfa}
\usepackage{xfrac}
\usepackage[unicode,hidelinks,bookmarks=false]{hyperref}
\newcommand{\ie}{i.e.\ }
\newcommand{\cf}{cf.\ }
\newcommand{\resp}{respectively\ }
\newcommand{\Subsection}{\subsection}
\providecommand{\nobreakdash}{\nobreak\hbox{-}}
\setlength{\emergencystretch}{2em}
\multlinegap0pt
\usepackage{bm}
\usepackage{bbm}
\usepackage{dsfont}
\usepackage{xspace}
\usepackage{cancel}
\usepackage{mathtools}
\usepackage{amsthm}
\makeatletter
\@ifundefined{definition}{\newtheorem{definition}{Definition}}{}
\@ifundefined{proposition}{\newtheorem{proposition}[definition]{Proposition}}{}
\@ifundefined{theorem}{\newtheorem{theorem}[definition]{Theorem}}{}
\makeatother
\newcommand{\HT}{\mathrm{HT}}
\newcommand{\ZZ}{{\mathbb Z}}
\newcommand{\cF}{{\mathcal F}}
\newcommand{\cO}{{\mathcal O}}
\newcommand{\frakX}{\mathfrak X}
\newcommand{\frako}{\mathfrak o}
\newcommand{\frakp}{\mathfrak p}
\newcommand{\tor}{\mathrm{tor}} %toroidal
\title[Eigenvariety for partially classical Hilbert modular forms]{Eigenvariety for partially classical Hilbert modular forms}
\author{Mladen Dimitrov, Chi-Yun Hsu}
\date{Source: arXiv:2403.09784v3 (CC BY 4.0)}
\begin{document}
\maketitle

\noindent\emph{Source and licence.} Eigenvariety for partially classical Hilbert modular forms. Version arXiv:2403.09784v3 (CC BY 4.0), distributed under the Creative Commons Attribution 4.0 International licence, as recorded in the arXiv licence metadata for this exact version and shown on the abstract page https://arxiv.org/abs/2403.09784v3. Full rights notice: licenses/arxiv-2403.09784-CC-BY-4.0.txt.\par\smallskip

\noindent\textit{Editorial scope.} Counted unit: the Theorem whose source label is \texttt{None} in the article (source0.tex lines 649--660, source label \texttt{None}), delivered together with the article's own complete original proof of it (source0.tex lines 662--681, one \texttt{proof} environment of 20 lines). No mathematics is written, repaired or re-derived: statement and proof are the article's own text line for line. Required context retained uncounted: prop:omega\_isom (lines 605--611); prop:cokerHT (lines 638--641). Every definition, display and label of those lines is retained unchanged. Omitted from this package and not redistributed: 1--604, 612--637, 642--648, 661--661, 682--1392. Nothing inside a retained excerpt is omitted or altered: every retained body is delivered exactly as the article's lines are written, so no omission receipt of any kind accompanies this package. Citations: the retained excerpts carry no citation command at all, so no bibliography entry of the article is transcribed and no reference is redistributed. Reference adaptation: none was needed and none was made; every reference inside the delivered unit resolves inside it, so no unresolved reference or citation is produced. Printed slips kept literal: no printed word, symbol, hypothesis or number of the counted statement or of its proof was altered. Layout and class: the article's own class is \texttt{amsart} with packages including fullpage, amsmath, amssymb, mathrsfs, bbm, tikz, amsthm, enumerate, enumitem, hyperref, colonequals; the wrapper is the lane's pinned \texttt{amsart} editorial wrapper at 10pt on A4, which restates the theorem-like environments the retained text needs and adds the article's own macro definitions that the retained text needs (\texttt{\textbackslash HT}, \texttt{\textbackslash ZZ}, \texttt{\textbackslash cF}, \texttt{\textbackslash cO}, \texttt{\textbackslash frakX}, \texttt{\textbackslash frako}, \texttt{\textbackslash frakp}, \texttt{\textbackslash tor}), with extra wrapper packages (bm, bbm, dsfont, xspace, cancel, mathtools). The delivered numbering is the unit's own. Assets: no retained span contains a figure environment, a graphics-inclusion command, a PDF-page inclusion, a table or a TikZ picture; no image file is copied into this package, the package declares no asset dependency and no image file is redistributed with it. Licence: version arXiv:2403.09784v3 of this article is distributed under the Creative Commons Attribution 4.0 International licence, as recorded in the arXiv licence metadata for this exact version and shown on the abstract page https://arxiv.org/abs/2403.09784v3. A full rights notice is delivered at licenses/arxiv-2403.09784-CC-BY-4.0.txt. Characters: every retained excerpt is pure ASCII, so the delivered engine, XeLaTeX with its default Latin Modern text font, reports no missing character.
\begin{proposition} \label{prop:omega_isom}
For $w\leqslant n-\frac{p^n-1}{p-1}v_\frakp$,
the morphism $\omega_{\frakp}  \rightarrow  \omega_{\frakp}/p^n\omega_{\frakp}
=\omega_{\widetilde{A}[\frakp^n]} \to  \omega_{\widetilde{C}_{\frakp,n}}$ of sheaves on $\frakX_0^\tor(p^{\underline{n}},\underline{v})$
induced by the inclusion $\widetilde{C}_{\frakp,n}\subseteq \widetilde{A}[\frakp^n]$
 becomes an isomorphisms modulo $p^w$.
\end{proposition}

\begin{proposition} \label{prop:cokerHT}
    The cokernel of $\HT_{\frakp,n}\otimes 1$
     is annihilated  by $p^\frac{v_\frakp}{p-1}$.
\end{proposition}

\begin{theorem}
There exists a sheaf $\cF_\frakp\subseteq \omega_\frakp$  containing $p^{\frac{v_\frakp}{p-1}}\omega_\frakp$ such that
\begin{enumerate}
    \item $\cF_\frakp$ is locally free of rank $f_\frakp$ on $\frakX_{1,\frakp}^\tor(p^{\underline{n}},\underline{v})$, and 
    \item  for all $w\leqslant n-\frac{p^n}{p-1}v_\frakp$ the  map
$\HT_{\frakp,w}\colon \widetilde{C}_{\frakp,n}^{D} \rightarrow \cF_{\frakp}/p^w\cF_{\frakp}$
deduced from $\HT_{\frakp,n}$ induces an isomorphism
\[
\HT_{\frakp,w} \otimes 1 \colon \widetilde{C}_{\frakp,n}^{D}\otimes \cO_{\frakX_{1,\frakp}^\tor(p^{\underline{n}},\underline{v})} \xrightarrow{\sim} \cF_\frakp/p^w\cF_{\frakp}.
\]
\end{enumerate}
\end{theorem}

\begin{proof} Let  $x_1,\ldots, x_{f_{\frakp}}$ be the image of a $\ZZ/p^n\ZZ$-basis of $\frako/\frakp^n$ 
under the universal trivialization $\psi_n\colon \frako/\frakp^n \rightarrow \widetilde{C}_{\frakp,n}^{D}$.
Let $y_1,\ldots, y_{f_\frakp}\in\omega_\frakp$ be lifts of $\HT_{\frakp,n}(x_i)\in \omega_{\widetilde{C}_{\frakp,n}}$ under the surjection $\omega_\frakp\twoheadrightarrow \omega_{\widetilde{C}_{\frakp,n}}$ and let $\cF_\frakp\subseteq \omega_\frakp$ be their linear span. 

We claim that $\cF_\frakp$ is locally free of rank $f_\frakp$. 
Indeed, suppose there is a non-trivial relation $\sum_{i=1}^{f_\frakp} a_i y_i = 0$.
Letting $w_0=n-\frac{p^n-1}{p-1}v_\frakp$, we may assume that there is at least one $a_i$, say $a_1$, not divisible by $p^{w_0-\frac{v_\frakp}{p-1}}$.
In particular, by Proposition~\ref{prop:omega_isom}, after projecting the relation to $\omega_{\widetilde{C}_{\frakp,n},w_0}$, we have $\sum_{i=1}^{f_\frakp} a_i \HT_{\frakp,n}(x_i) = 0$.
By Proposition~\ref{prop:cokerHT}, the kernel of $\HT_{\frakp,n}$ is killed by $p^{\frac{v_\frakp}{p-1}}$, and so $\sum_{i=1}^{f_\frakp} p^{\frac{v_\frakp}{p-1}}a_i x_i=0$. 
However, $p^{\frac{v_\frakp}{p-1}}a_1$ is not divisible by $p^{w_0}$ and hence not divisible by $p^n$, contradicting the fact that the $x_i$'s are images under $\psi_n$ of a basis of $\frako/\frakp^n$.

By Proposition~\ref{prop:cokerHT}, $\cF_\frakp$ contains $p^{\frac{v_\frakp}{p-1}}\omega_\frakp$.
In addition, $\cF_\frakp$ is independent of the choice of the lifts $y_i$.
This is because by the proof of Proposition~\ref{prop:omega_isom}, the kernel of $\omega_\frakp\twoheadrightarrow \omega_{\widetilde{C}_{\frakp.n}}$ is contained in $p^{n-\frac{p^n-1}{p-1}v_\frakp}\omega_\frakp\subseteq p^{\frac{v_\frakp}{p-1}}\omega_\frakp$ and hence contained in $\cF_\frakp$.

Composing $\HT_{\frakp,n}\colon \widetilde{C}_{\frakp,n}^{D} \rightarrow \omega_{\widetilde{C}_{\frakp,n}}$ with reduction modulo $p^{w_0}$ and using Proposition~\ref{prop:omega_isom}, we have the map $\widetilde{C}_{\frakp,n}^{D} \rightarrow \omega_{\frakp}/p^w\omega_\frakp$, which factors through $\cF_\frakp/(p^{w_0}\omega_\frakp \cap \cF_\frakp)$ by construction of $\cF_\frakp$.
Let $w\leqslant n-\frac{p^n}{p-1}v_\frakp$.
Then $p^{w_0}\omega_\frakp \cap \cF_\frakp \subseteq p^w \cF_\frakp$ because $\cF_\frakp$ contains $p^{\frac{v_\frakp}{p-1}}\omega_\frakp$.
Hence we have $\HT_{\frakp,w}\colon \widetilde{C}_{\frakp,n}^{D} \rightarrow \cF_{\frakp}/p^w \cF_{\frakp}$.
Finally, $\HT_{\frakp,w}\otimes 1$ is an isomorphism because it is a surjection between two locally free sheaves of the same rank $f_\frakp$.\end{proof}

\end{document}
